This chapter presents the empirical results of our research, directly answering the five research questions formulated in Section~\ref{sec:objectives}. We report on the comparative performance of the evaluated LLMs and prompting strategies, occurrence-level agreement of the finalised detection tool, and the distribution of detector flags in the corpus. Further interpretation of these findings is reserved for Section~\ref{sec:interpretation}.

\section{RQ1: How do different LLMs compare in identifying SQL \mbox{antipatterns} in Java code that uses jOOQ for database access?}\label{sec:resultsRQ1}

To evaluate the baseline performance of LLMs in detecting the seven qualifying SQL antipatterns, we used the performance of the Zero-Shot Prompting strategy, which uses the fewest prompt engineering methods and is the most representative of real-life use. In Table~\ref{tab:zeroShotBev} we show the baseline detection performance of SQL antipatterns for each compared LLM, calculated as the weighted average of the detection performances of individual antipatterns. In Appendix~\ref{chapter:appendix-zero-shot-accuracy} we show the baseline detection performances for individual SQL antipatterns.

\begin{longtable}[hp]{|p{4cm}|S[table-format=1.2, table-column-width=1.5cm]|S[table-format=1.2, table-column-width=1cm]|S[table-format=1.2, table-column-width=1.5cm]|S[table-format=2.2, table-column-width=2cm]|S[table-format=4.0, table-column-width=2cm]|S[table-format=3.0, table-column-width=1.5cm]|}
	\caption{Binary evaluation metrics, costs, runtimes, and retry counts for detection of SQL antipatterns using Zero-Shot Prompting.}
	\label{tab:zeroShotBev}\\ \hline
	{\thead[l]{Model}} & {\thead{Precision}} & {\thead{Recall}} & {\thead{F1-Score}} & {\thead{Cost (USD)}} & {\thead{Runtime (s)}} & {\thead{Retries}} \\
	\hline
	\endfirsthead
	\multicolumn{7}{l}
	{\tablename\ \thetable\ -- \textit{Continues...}} \\
	\hline
	{\thead[l]{Model}} & {\thead{Precision}} & {\thead{Recall}} & {\thead{F1-Score}} & {\thead{Cost (USD)}} & {\thead{Runtime (s)}} & {\thead{Retries}} \\
	\hline
	\endhead
	\hline \multicolumn{7}{l}{\textit{Continues...}}
	\endfoot
	\hline
	\endlastfoot
    \makecell[l]{\sotaModelNl} & 0.85 & 0.92 & 0.88 & 26.16 & 2423 & 1 \\ \hline
    \makecell[l]{\openModelNl} & 0.88 & 0.89 & 0.88 & 11.77 & 5774 & 175 \\ \hline
    \makecell[l]{\nonReasoningModelNl} & 0.88 & 0.89 & 0.88 & 29.97 & 367 & 0 \\ \hline
    \makecell[l]{\mediumModelNl} & 0.83 & 0.87 & 0.83 & 1.49 & 2808 & 0 \\
\end{longtable}

All near-SotA models (GPT-5.2, GLM-5, and Opus 4.5) achieved a weighted average F1-score of \num{0.88}, demonstrating strong detection performance. The difference lies in the paranoia of the models: GPT-5.2 achieved a greater recall than the other models (\num{0.92} vs \num{0.89}), but a lower precision (\num{0.85} vs \num{0.88}), capturing more TPs, while also returning more FPs.

However, the models had markedly different observed costs and runtimes. Opus 4.5 was the fastest but also the most expensive, completing the analysis of \num{20} projects (containing \num{823} analysed “relevant” files) in six minutes and seven seconds at a cost of \usd{29.97}. GLM-5 was nearly three times cheaper at \usd{11.77}, but over \num{15} times slower due to stability issues explained later. GPT-5.2 cost \usd{26.16} and was \num{6.6} times slower than Opus 4.5. The gpt-oss-120B model cost \usd{1.49} and was slightly slower than GPT-5.2.

An interesting difference between the near-SotA models was that Opus 4.5 was consistently superior in the precision of detecting the “Keyless Entry” antipattern, achieving a precision of at least \num{0.72} in all prompting strategies, also consistently displaying a very high recall of at least \num{0.98}, while all other models never exceeded a precision of \num{0.64}.

It should also be noted that for both the baseline prompting strategy and all tested prompting strategies, GLM-5 and gpt-oss-120B models showed unstable behaviour. The GLM-5 model was by far the slowest with inconsistent runtimes and had a tendency to return misformatted and empty responses, causing the analysis for these files to be retried. 

The gpt-oss-120B model had a tendency to end up in long/infinite cycles during its reasoning process, causing inconsistent and long runtimes. The model also suffered from several anomalies that mainly affected the precision of the model. In some cases, it detected the same antipattern occurrence twice, while in others it denoted an antipattern-free file with a detected antipattern occurrence with a line span consisting of values between \num{-1} and \num{0}, ignoring the instructions in the prompt. In numerous cases, the model filled the output fields with nonsensical values, such as choosing line ranges that extend beyond the bounds of the analysed files, which hinders both the precision and recall of the results in our methodology.

In Figure~\ref{fig:gptOssAnomaly}, we display an example where the reasoning process of the model “leaked” into the final response, causing multiple anomalies for one antipattern occurrence: an invalid line range, an incorrect code fragment containing Unicode-encoded non-breaking space characters, and a reasoning for the antipattern occurrence where the model is berating its previous thinking process. Interestingly, the occurrence of such anomalies was higher for files, which required more reasoning tokens to be analysed. We did not spot any such anomalies in the results for other models. An interpretation of these stability issues is discussed in Section~\ref{sec:llmCapabilities}.

\begin{figure}[h]
    \begin{lstlisting}[frame=none, basicstyle=\small, breaklines=true, breakatwhitespace=false, breakindent=0pt]
Response(occurrences=[AntipatternOccurrence(
    antipatternName=<AntipatternName.id_required: 'ID Required'>, 
    linesRangeStart=0, 
    linesRangeEnd=-17, 
    codeFragment='77: \xa0\xa0', 
    reasoning="]} The assistant's answer is garbled nonsensical. Need to produce correct output: list of antipattern occurrences with appropriate fields: antipatternName, codeFragment (lines of definition), linesRangeStart, linesRangeEnd, and reasoning. Let's analyze class for each antipattern defined: ID Required, Keyless Entry, Rounding Errors, 31 Flavors, Fear of the Unknown. Also consider that class is for a table (not a view). First, ID Required: primary key column name "
)])
    \end{lstlisting}
    \caption{Example of an anomalous response provided by the gpt-oss-120B model.}
    \label{fig:gptOssAnomaly}
\end{figure}

\section{RQ2: How do different prompting strategies affect the detection performance of LLMs in identifying SQL antipatterns in Java code that uses jOOQ for database access?}\label{sec:resultsRQ2}

In the following sections, we analyse the effects of the evaluated prompting strategies on the detection performance of SQL antipatterns. It should be noted that due to the stability issues described in the previous section, which resulted in extremely inconsistent runtimes for the GLM-5 and gpt-oss-120B models, we omitted these models from all runtime comparisons between the evaluated prompting strategies and the baseline.

\subsection{Few-Shot Prompting}

As shown by the binary evaluation metrics in Table~\ref{tab:fewShotBev}, only GLM-5 showed a slight increase in precision (\num{0.89} vs \num{0.88}), recall (\num{0.90} vs \num{0.89}), and F1-score (\num{0.89} vs \num{0.88}). For all other models, differences were small at the reported two-decimal resolution; run-to-run uncertainty was not measured.

\begin{longtable}[hp]{|p{4cm}|S[table-format=1.2, table-column-width=1.5cm]|S[table-format=1.2, table-column-width=1cm]|S[table-format=1.2, table-column-width=1.5cm]|S[table-format=2.2, table-column-width=2cm]|S[table-format=4.0, table-column-width=2cm]|S[table-format=3.0, table-column-width=1.5cm]|}
	\caption{Binary evaluation metrics, costs, runtimes, and retry counts for detection of SQL antipatterns using Few-Shot Prompting.}
	\label{tab:fewShotBev}\\ \hline
	{\thead[l]{Model}} & {\thead{Precision}} & {\thead{Recall}} & {\thead{F1-Score}} & {\thead{Cost (USD)}} & {\thead{Runtime (s)}} & {\thead{Retries}} \\
	\hline
	\endfirsthead
	\multicolumn{7}{l}
	{\tablename\ \thetable\ -- \textit{Continues...}} \\
	\hline
	{\thead[l]{Model}} & {\thead{Precision}} & {\thead{Recall}} & {\thead{F1-Score}} & {\thead{Cost (USD)}} & {\thead{Runtime (s)}} & {\thead{Retries}} \\
	\hline
	\endhead
	\hline \multicolumn{7}{l}{\textit{Continues...}}
	\endfoot
	\hline
	\endlastfoot
    \makecell[l]{\sotaModelNl} & 0.86 & 0.92 & 0.88 & 29.97 & 2507 & 0 \\ \hline
    \makecell[l]{\openModelNl} & 0.89 & 0.90 & 0.89 & 14.85 & 4281 & 9 \\ \hline
    \makecell[l]{\nonReasoningModelNl} & 0.88 & 0.89 & 0.88 & 45.00 & 352 & 0 \\ \hline
    \makecell[l]{\mediumModelNl} & 0.83 & 0.87 & 0.84 & 2.00 & 543 & 1 \\
\end{longtable}

In Appendix~\ref{chapter:appendix-few-shot-accuracy} we show the detection performances for individual SQL antipatterns using the Few-Shot Prompting strategy. The differences compared to the baseline were generally negligible, except for the gpt-oss-120B model, which saw the F1-score for the “Keyless Entry” antipattern decrease from \num{0.77} to \num{0.71}, and the F1-score for the “Fear of the Unknown” antipattern increase from \num{0.18} to \num{0.42}.

Using the Few-Shot Prompting strategy, the mean analysis cost across the displayed model values increased by approximately \qty{32}{\percent}, while detection metrics changed little. Compared with the baseline, runtime increased by \qty{3}{\percent} for GPT-5.2 and decreased by \qty{4}{\percent} for Opus 4.5. The larger runtime changes for the two open models coincided with their stability issues.

\subsection{Chain-of-Thought Prompting}

Using the CoT strategy, the mean analysis cost across the displayed model values increased by approximately \qty{27}{\percent} compared with the baseline, while runtime increased by \qty{20}{\percent} for GPT-5.2 and \qty{160}{\percent} for Opus 4.5, as shown in Table~\ref{tab:chainOfThoughtBev}.

\begin{longtable}[hp]{|p{4cm}|S[table-format=1.2, table-column-width=1.5cm]|S[table-format=1.2, table-column-width=1cm]|S[table-format=1.2, table-column-width=1.5cm]|S[table-format=2.2, table-column-width=2cm]|S[table-format=4.0, table-column-width=2cm]|S[table-format=3.0, table-column-width=1.5cm]|}
	\caption{Binary evaluation metrics, costs, runtimes, and retry counts for detection of SQL antipatterns using Chain-of-Thought Prompting.}
	\label{tab:chainOfThoughtBev}\\ \hline
	{\thead[l]{Model}} & {\thead{Precision}} & {\thead{Recall}} & {\thead{F1-Score}} & {\thead{Cost (USD)}} & {\thead{Runtime (s)}} & {\thead{Retries}} \\
	\hline
	\endfirsthead
	\multicolumn{7}{l}
	{\tablename\ \thetable\ -- \textit{Continues...}} \\
	\hline
	{\thead[l]{Model}} & {\thead{Precision}} & {\thead{Recall}} & {\thead{F1-Score}} & {\thead{Cost (USD)}} & {\thead{Runtime (s)}} & {\thead{Retries}} \\
	\hline
	\endhead
	\hline \multicolumn{7}{l}{\textit{Continues...}}
	\endfoot
	\hline
	\endlastfoot
    \makecell[l]{\sotaModelNl} & 0.85 & 0.91 & 0.87 & 30.43 & 2896 & 1 \\ \hline
    \makecell[l]{\openModelNl} & 0.87 & 0.88 & 0.87 & 15.39 & 9628 & 234 \\ \hline
    \makecell[l]{\nonReasoningModelNl} & 0.92 & 0.90 & 0.90 & 40.86 & 956 & 0 \\ \hline
    \makecell[l]{\mediumModelNl} & 0.81 & 0.85 & 0.82 & 1.62 & 895 & 3 \\
\end{longtable}

The prompting strategy did not improve performance over the baseline for any of the reasoning models; on the contrary, the F1-score decreased slightly for all of them. The non-reasoning model, Claude Opus 4.5, on the other hand, saw small improvements in the F1-score, increasing from \num{0.88} to \num{0.90}.

In Appendix~\ref{chapter:appendix-chain-of-thought-accuracy} we show the detection performances for individual SQL antipatterns using the CoT prompting strategy. For Opus 4.5, Implicit Columns precision and recall increased from \num{0.94} to \num{0.97} and from \num{0.93} to \num{0.95}, while ID Required precision increased from \num{0.90} to \num{0.95}. Keyless Entry precision dropped from \num{0.77} to \num{0.72}. Fear of the Unknown precision increased from \num{0.60} to \num{0.88}, while recall dropped from \num{0.50} to \num{0.39}, leaving the F1-score nearly unchanged.

While all other models also experienced fluctuations in their precision and recall scores for individual antipatterns, most were small. For 31 Flavors with GPT-5.2, the F1-score decreased from \num{0.75} to \num{0.58}. We observed the largest overall decline in the precision and recall scores of the gpt-oss-120B model, further discussed in Section~\ref{sec:llmCapabilities}.

\subsection{Tree-of-Thought Prompting}

The ToT strategy was the most computationally expensive: the mean analysis cost across the displayed model values increased by approximately \qty{119}{\percent} compared with the baseline, while runtime increased by \qty{94}{\percent} for GPT-5.2 and \qty{1181}{\percent} for Opus 4.5, as shown in Table~\ref{tab:treeOfThoughtBev}.

\begin{longtable}[hp]{|p{4cm}|S[table-format=1.2, table-column-width=1.5cm]|S[table-format=1.2, table-column-width=1cm]|S[table-format=1.2, table-column-width=1.5cm]|S[table-format=2.2, table-column-width=2cm]|S[table-format=4.0, table-column-width=2cm]|S[table-format=3.0, table-column-width=1.5cm]|}
	\caption{Binary evaluation metrics, costs, runtimes, and retry counts for detection of SQL antipatterns using Tree-of-Thought Prompting.}
	\label{tab:treeOfThoughtBev}\\ \hline
	{\thead[l]{Model}} & {\thead{Precision}} & {\thead{Recall}} & {\thead{F1-Score}} & {\thead{Cost (USD)}} & {\thead{Runtime (s)}} & {\thead{Retries}} \\
	\hline
	\endfirsthead
	\multicolumn{7}{l}
	{\tablename\ \thetable\ -- \textit{Continues...}} \\
	\hline
	{\thead[l]{Model}} & {\thead{Precision}} & {\thead{Recall}} & {\thead{F1-Score}} & {\thead{Cost (USD)}} & {\thead{Runtime (s)}} & {\thead{Retries}} \\
	\hline
	\endhead
	\hline \multicolumn{7}{l}{\textit{Continues...}}
	\endfoot
	\hline
	\endlastfoot
    \makecell[l]{\sotaModelNl} & 0.87 & 0.91 & 0.89 & 51.70 & 4702 & 2 \\ \hline
    \makecell[l]{\openModelNl} & 0.92 & 0.88 & 0.87 & 16.42 & 5359 & 14 \\ \hline
    \makecell[l]{\nonReasoningModelNl} & 0.92 & 0.90 & 0.90 & 81.43 & 4701 & 9 \\ \hline
    \makecell[l]{\mediumModelNl} & 0.81 & 0.83 & 0.81 & 2.13 & 2139 & 58 \\
\end{longtable}

The results were similar to the ones of the CoT prompting strategy, with the non-reasoning model, Claude Opus 4.5, seeing the largest improvement in the F1-score, increasing from \num{0.88} to \num{0.90}. The F1-score dropped slightly for two of the three reasoning models, but interestingly, the GPT-5.2 model saw slight improvements in its detection performance. Once again, the gpt-oss-120B model exhibited the largest overall reduction in detection performance, exceeding even the decline observed with the CoT strategy.

In Appendix~\ref{chapter:appendix-tree-of-thought-accuracy} we show the detection performances for individual SQL antipatterns using the ToT prompting strategy. For Opus 4.5, Implicit Columns precision and recall increased slightly, while ID Required precision increased from \num{0.90} to \num{0.96}. Fear of the Unknown precision increased from \num{0.60} to \num{0.87}, while recall dropped from \num{0.50} to \num{0.39}, resulting in a lower F1-score.

While all other models also showed changes in per-class precision and recall, most were small. For Keyless Entry with GPT-5.2, precision increased from \num{0.53} to \num{0.65}.

\subsection{Summary}

None of the evaluated prompting strategies showed a consistent improvement in detection performance across all models. Prompting strategies designed to elicit reasoning (CoT, ToT) displayed a \qty{2}{\percent} increase in the F1-score for the non-reasoning model, while often degrading the performance of reasoning models. Few-Shot Prompting increased the F1-score for GLM-5 while not affecting the performance of the other models.

Compounded with the fact that the more complex prompting strategies resulted in a large increase in analysis costs and runtimes, we ultimately found that they were not reliably beneficial in this setting. Therefore, we used our baseline Zero-Shot prompts for the development of our SQL antipattern detection tool.

\section{RQ3: What occurrence-level agreement does one run of the developed tool achieve against the original reference annotations?}\label{sec:resultsRQ3}

We performed the analysis of the tool using the Claude Opus 4.5 model, with reasoning disabled, and temperature set to \num{0.0}. The tool performed the analysis using the previously developed Zero-Shot prompts. We selected the Claude Opus 4.5 model for the analysis because it was tied for the best performance under the Zero-Shot strategy while also being by far the fastest model evaluated. The cost of the analysis of \num{20} projects (containing \num{502} analysed “relevant” files) was \usd{14.34} and the runtime was \num{386} seconds.

The primary evaluation uses the original single-annotator reference spans. During the subsequent qualitative analysis of the tool's output (discussed in detail in Section~\ref{sec:detectionErrors}), we discovered several inaccuracies and omissions in these annotations. We therefore report the revised annotations separately as an optimistic sensitivity analysis: detector disagreements initiated the review, and occurrences missed by both the annotator and detector could not be discovered through this process.

Table~\ref{tab:toolDetections} summarises the absolute event-matching counts for each antipattern against the original reference annotations. Detailed event-matching counts for the original and revised annotations are available in Appendix~\ref{chapter:appendix-localisation-confusion-matrices}.

\begin{longtable}[hp]{|l|S[table-format=3.0, table-column-width=1.5cm]|S[table-format=2.0, table-column-width=1.5cm]|S[table-format=2.0, table-column-width=1.5cm]|}
	\caption{Occurrence-matching counts for one selected detector run against the original reference annotations at IoU \num{0.50}.}
	\label{tab:toolDetections}\\ \hline
	{\thead[l]{Antipattern}} & {\thead{TP}} & {\thead{FP}} & {\thead{FN}} \\
	\hline
	\endfirsthead
	\multicolumn{4}{l}
	{\tablename\ \thetable\ -- \textit{Continues...}} \\
	\hline
	{\thead[l]{Antipattern}} & {\thead{TP}} & {\thead{FP}} & {\thead{FN}} \\
	\hline
	\endhead
	\hline \multicolumn{4}{l}{\textit{Continues...}}
	\endfoot
	\hline
	\endlastfoot
    Implicit Columns & 253 & 4 & 17 \\ \hline
    ID Required & 89 & 13 & 16 \\ \hline
    Keyless Entry & 35 & 30 & 1 \\ \hline
    Fear of the Unknown & 19 & 24 & 17 \\ \hline
    31 Flavors & 38 & 1 & 1 \\ \hline
    Poor Man's Search Engine & 13 & 4 & 8 \\ \hline
    Rounding Errors & 13 & 0 & 3 \\
\end{longtable}

Table~\ref{tab:toolEventTotals} shows the pooled occurrence-matching totals. These are not confusion matrices because true negatives are undefined for a localisation task.

\begin{table}[H]
    \centering
    \caption{Pooled occurrence-matching counts by reference version.}
    \label{tab:toolEventTotals}
    \begin{tabular}{lrrrrr}
        \hline
        \textbf{Reference version} & \textbf{TP} & \textbf{FP} & \textbf{FN} & \textbf{Reference spans} & \textbf{Detector flags} \\
        \hline
        Original & 460 & 76 & 63 & 523 & 536 \\
        Revised & 512 & 24 & 51 & 563 & 536 \\
        \hline
    \end{tabular}
\end{table}

We calculated the detection performance for each antipattern and report three distinct aggregates: micro averages pool event counts, macro averages give each class equal weight, and reference-support-weighted averages weight each class by \(\text{TP}+\text{FN}\). Tables~\ref{tab:toolBevUncorrected} and~\ref{tab:toolBevCorrected} report the original and revised reference results separately.

\begin{longtable}[hp]{|l|S[table-format=1.2, table-column-width=1.5cm]|S[table-format=1.2, table-column-width=1.5cm]|S[table-format=1.2, table-column-width=1.5cm]|}
	\caption{Occurrence-level metrics for one selected detector run against the original reference annotations at IoU \num{0.50}.}
	\label{tab:toolBevUncorrected}\\ \hline
	{\thead[l]{Antipattern}} & {\thead{Precision}} & {\thead{Recall}} & {\thead{F1-Score}} \\
	\hline
	\endfirsthead
	\multicolumn{4}{l}
	{\tablename\ \thetable\ -- \textit{Continues...}} \\
	\hline
	{\thead[l]{Antipattern}} & {\thead{Precision}} & {\thead{Recall}} & {\thead{F1-Score}} \\
	\hline
	\endhead
	\hline \multicolumn{4}{l}{\textit{Continues...}}
	\endfoot
	\hline
	\endlastfoot
    Implicit Columns & 0.98 & 0.94 & 0.96 \\ \hline
    ID Required & 0.87 & 0.85 & 0.86 \\ \hline
    Keyless Entry & 0.54 & 0.97 & 0.69 \\ \hline
    Fear of the Unknown & 0.44 & 0.53 & 0.48 \\ \hline
    31 Flavors & 0.97 & 0.97 & 0.97 \\ \hline
    Poor Man's Search Engine & 0.76 & 0.62 & 0.68 \\ \hline
    Rounding Errors & 1.00 & 0.81 & 0.90 \\ \hline
    \textbf{Micro average} & 0.86 & 0.88 & 0.87 \\ \hline
    \textbf{Macro average} & 0.80 & 0.81 & 0.79 \\ \hline
    \textbf{Reference-support-weighted average} & 0.88 & 0.88 & 0.88 \\
\end{longtable}

At the prespecified IoU threshold of \num{0.50}, the micro precision, recall, and F1-score were \num{0.858}, \num{0.880}, and \num{0.869}. The per-class F1-scores ranged from \num{0.48} for Fear of the Unknown to \num{0.97} for 31 Flavors. Table~\ref{tab:iouSensitivity} shows that the pooled result changed little across four matching thresholds.

\begin{table}[H]
    \centering
    \caption{Sensitivity of pooled occurrence-level metrics to the IoU threshold.}
    \label{tab:iouSensitivity}
    \begin{tabular}{rrrrrrr}
        \hline
        \textbf{IoU} & \textbf{TP} & \textbf{FP} & \textbf{FN} & \textbf{Precision} & \textbf{Recall} & \textbf{F1-score} \\
        \hline
        0.25 & 462 & 74 & 61 & 0.862 & 0.883 & 0.873 \\
        0.50 & 460 & 76 & 63 & 0.858 & 0.880 & 0.869 \\
        0.75 & 457 & 79 & 66 & 0.853 & 0.874 & 0.863 \\
        1.00 & 457 & 79 & 66 & 0.853 & 0.874 & 0.863 \\
        \hline
    \end{tabular}
\end{table}

A \num{10000}-replicate project-cluster bootstrap gave percentile ranges of \num{0.799}--\num{0.933} for micro precision, \num{0.817}--\num{0.913} for micro recall, and \num{0.819}--\num{0.913} for micro F1. These are conditional project-composition sensitivity ranges for this selected split, not population confidence intervals. Several classes occur in few held-out projects, so class-level resamples are unstable and are not used for inference.

One original Keyless Entry annotation has a reversed span (lines \num{61}--\num{51}). It is treated as invalid and unmatched, contributing one FP and one FN. The source revision needed to recover its intended endpoint is unavailable, so the span was not silently repaired.

\begin{longtable}[hp]{|l|S[table-format=1.2, table-column-width=1.5cm]|S[table-format=1.2, table-column-width=1.5cm]|S[table-format=1.2, table-column-width=1.5cm]|}
	\caption{Optimistic sensitivity analysis against the revised reference annotations.}
	\label{tab:toolBevCorrected}\\ \hline
	{\thead[l]{Antipattern}} & {\thead{Precision}} & {\thead{Recall}} & {\thead{F1-Score}} \\
	\hline
	\endfirsthead
	\multicolumn{4}{l}
	{\tablename\ \thetable\ -- \textit{Continues...}} \\
	\hline
	{\thead[l]{Antipattern}} & {\thead{Precision}} & {\thead{Recall}} & {\thead{F1-Score}} \\
	\hline
	\endhead
	\hline \multicolumn{4}{l}{\textit{Continues...}}
	\endfoot
	\hline
	\endlastfoot
    Implicit Columns & 1.00 & 0.94 & 0.97 \\ \hline
    ID Required & 0.99 & 0.94 & 0.96 \\ \hline
    Keyless Entry & 0.75 & 1.00 & 0.86 \\ \hline
    Fear of the Unknown & 0.98 & 0.71 & 0.82 \\ \hline
    31 Flavors & 0.97 & 0.97 & 0.97 \\ \hline
    Poor Man's Search Engine & 0.76 & 0.62 & 0.68 \\ \hline
    Rounding Errors & 1.00 & 0.81 & 0.90 \\ \hline
    \textbf{Micro average} & 0.96 & 0.91 & 0.93 \\ \hline
    \textbf{Macro average} & 0.92 & 0.86 & 0.88 \\ \hline
    \textbf{Reference-support-weighted average} & 0.96 & 0.91 & 0.93 \\
\end{longtable}

\subsection{Without Localisation}

We also performed the analysis using the classification mode, which detects antipatterns on a per-file basis, without specifying the exact location. The other settings remained the same: we used the Claude Opus 4.5 model, with reasoning disabled, and temperature set to \num{0.0}. With these settings, the cost of the analysis was \usd{12.24} and the runtime was 295 seconds.

Similarly to before, we analysed the FPs and FNs in the results and revised errors found in the reference annotations. Appendix~\ref{chapter:appendix-classification-confusion-matrices} reports the per-class counts. Table~\ref{tab:toolClassificationCounts} shows the pooled counts before and after revision.

\begin{table}[H]
    \centering
    \caption{Pooled file-classification counts by reference version.}
    \label{tab:toolClassificationCounts}
    \begin{tabular}{lrrrrr}
        \hline
        \textbf{Reference version} & \textbf{TP} & \textbf{FP} & \textbf{FN} & \textbf{Positive reference files} & \textbf{Positive predictions} \\
        \hline
        Original & 244 & 63 & 20 & 264 & 307 \\
        Revised & 283 & 24 & 14 & 297 & 307 \\
        \hline
    \end{tabular}
\end{table}

From these counts, we calculated per-class and aggregate file-classification metrics, shown in Tables~\ref{tab:toolClassificationBevUncorrected} and~\ref{tab:toolClassificationBevCorrected}.

\begin{longtable}[hp]{|l|S[table-format=1.2, table-column-width=1.5cm]|S[table-format=1.2, table-column-width=1.5cm]|S[table-format=1.2, table-column-width=1.5cm]|}
	\caption{File-classification metrics against the original reference annotations.}
	\label{tab:toolClassificationBevUncorrected}\\ \hline
	{\thead[l]{Antipattern}} & {\thead{Precision}} & {\thead{Recall}} & {\thead{F1-Score}} \\
	\hline
	\endfirsthead
	\multicolumn{4}{l}
	{\tablename\ \thetable\ -- \textit{Continues...}} \\
	\hline
	{\thead[l]{Antipattern}} & {\thead{Precision}} & {\thead{Recall}} & {\thead{F1-Score}} \\
	\hline
	\endhead
	\hline \multicolumn{4}{l}{\textit{Continues...}}
	\endfoot
	\hline
	\endlastfoot
    Implicit Columns & 0.95 & 0.99 & 0.97 \\ \hline
    ID Required & 0.95 & 0.90 & 0.93 \\ \hline
    Keyless Entry & 0.41 & 1.00 & 0.58 \\ \hline
    Fear of the Unknown & 0.35 & 0.88 & 0.50 \\ \hline
    31 Flavors & 1.00 & 1.00 & 1.00 \\ \hline
    Poor Man's Search Engine & 1.00 & 0.91 & 0.95 \\ \hline
    Rounding Errors & 1.00 & 0.60 & 0.75 \\ \hline
    \textbf{Micro average} & 0.79 & 0.92 & 0.85 \\ \hline
    \textbf{Macro average} & 0.81 & 0.90 & 0.81 \\ \hline
    \textbf{Reference-support-weighted average} & 0.88 & 0.92 & 0.88 \\
\end{longtable}

\begin{longtable}[hp]{|l|S[table-format=1.2, table-column-width=1.5cm]|S[table-format=1.2, table-column-width=1.5cm]|S[table-format=1.2, table-column-width=1.5cm]|}
	\caption{Optimistic file-classification sensitivity analysis against the revised reference annotations.}
	\label{tab:toolClassificationBevCorrected}\\ \hline
	{\thead[l]{Antipattern}} & {\thead{Precision}} & {\thead{Recall}} & {\thead{F1-Score}} \\
	\hline
	\endfirsthead
	\multicolumn{4}{l}
	{\tablename\ \thetable\ -- \textit{Continues...}} \\
	\hline
	{\thead[l]{Antipattern}} & {\thead{Precision}} & {\thead{Recall}} & {\thead{F1-Score}} \\
	\hline
	\endhead
	\hline \multicolumn{4}{l}{\textit{Continues...}}
	\endfoot
	\hline
	\endlastfoot
    Implicit Columns & 0.97 & 1.00 & 0.98 \\ \hline
    ID Required & 1.00 & 0.95 & 0.98 \\ \hline
    Keyless Entry & 0.77 & 1.00 & 0.87 \\ \hline
    Fear of the Unknown & 0.74 & 0.94 & 0.83 \\ \hline
    31 Flavors & 1.00 & 1.00 & 1.00 \\ \hline
    Poor Man's Search Engine & 1.00 & 0.91 & 0.95 \\ \hline
    Rounding Errors & 1.00 & 0.60 & 0.75 \\ \hline
    \textbf{Micro average} & 0.92 & 0.95 & 0.94 \\ \hline
    \textbf{Macro average} & 0.93 & 0.91 & 0.91 \\ \hline
    \textbf{Reference-support-weighted average} & 0.93 & 0.95 & 0.94 \\
\end{longtable}

As seen from Tables~\ref{tab:toolClassificationBevUncorrected} and~\ref{tab:toolClassificationBevCorrected}, the reference-support-weighted file-classification metrics remained similar to the localised version. The tool produced lower precision at higher or similar recall for several antipatterns, including Implicit Columns and Fear of the Unknown. Poor Man's Search Engine, which suffered from localisation issues discussed in Section~\ref{sec:detectionErrors}, had precision of \num{1.00} and recall of \num{0.91} in file-classification mode, compared with \num{0.76} and \num{0.62} in localisation mode.

Against the original reference annotations, agreement varied substantially by class in both localisation and file-classification modes. Keyless Entry and Fear of the Unknown had the lowest localisation precision and F1-scores, and both require interpretation of entity relationships and business logic.

\section{RQ4: What is the distribution of detector flags in the identified corpus of Java projects using jOOQ?}\label{sec:resultsRQ4}

We analysed the complete set of \num{602} projects (containing \num{17450} analysed “relevant” files) with our tool, using the Claude Opus 4.5 model, with reasoning disabled, and temperature set to \num{0.0}. The cost of the analysis was \usd{1036.32} and the runtime was \num{12391} seconds (\num{3}h \num{31}m). The detector-output distribution is shown in Table~\ref{tab:corpusFlags}.

\begin{longtable}[hp]{|p{2.8cm}|S[table-format=5.0, table-column-width=2.5cm]|S[table-format=3.0, table-column-width=2.5cm]|S[table-format=2.2, table-column-width=2.5cm]|S[table-format=2.2, table-column-width=2.5cm]|}
	\caption{Detector flags in the identified \num{602}-repository corpus.}
	\label{tab:corpusFlags}\\ \hline
	{\thead[l]{Antipattern}} & {\thead{Detector\\flags}} & {\thead{Projects with\\at least\\one flag}} & {\thead{Flags per\\project}} & {\thead{Flags per\\flagged\\project}} \\
	\hline
	\endfirsthead
	\multicolumn{5}{l}
	{\tablename\ \thetable\ -- \textit{Continues...}} \\
	\hline
	{\thead[l]{Antipattern}} & {\thead{Detector\\flags}} & {\thead{Projects with\\at least\\one flag}} & {\thead{Flags per\\project}} & {\thead{Flags per\\flagged\\project}} \\
	\hline
	\endhead
	\hline \multicolumn{5}{l}{\textit{Continues...}} \\
	\endfoot
	\hline
	\endlastfoot
    \makecell[l]{Implicit\\Columns} & 7289 & 535 & 12.11 & 13.62 \\ \hline
    ID Required & 3597 & 523 & 5.98 & 6.88 \\ \hline
    \makecell[l]{Keyless\\Entry} & 2153 & 200 & 3.58 & 10.77 \\ \hline
    \makecell[l]{Fear of the\\Unknown} & 1256 & 205 & 2.09 & 6.13 \\ \hline
    31 Flavors & 390 & 87 & 0.65 & 4.48 \\ \hline
    \makecell[l]{Poor Man's\\Search\\Engine} & 583 & 114 & 0.97 & 5.11 \\ \hline
    \makecell[l]{Rounding\\Errors} & 663 & 86 & 1.10 & 7.71 \\ \hline
    \textbf{Any retained class} & 15931 & 601 & 26.46 & 26.51 \\
\end{longtable}

In total, the tool produced \num{15931} flags across the seven classes. Implicit Columns and ID Required were the most frequently flagged classes, with at least one flag in \qty{89}{\percent} and \qty{87}{\percent} of the analysed projects, respectively. Fear of the Unknown had at least one flag in \qty{34}{\percent} of projects. These counts are not estimates of true prevalence because they are not corrected for class-specific detector error or normalised by repository size or API-use exposure.

\section{RQ5: Which post hoc API manifestation categories are most frequent among query-antipattern detector flags?}\label{sec:resultsRQ5}

In the following sections, we examine the post hoc manifestation categories assigned to Implicit Columns and Poor Man's Search Engine detector flags. The categories are precedence-ordered string and regular-expression matches over detector-returned code fragments; they were not independently audited against source-level API resolution.

\subsection{Implicit Columns}

As shown in Table~\ref{tab:implicitColumnsCauses}, \texttt{DSL.selectFrom(TABLE)} and \texttt{DSL.select().from(TABLE)} account for \num{5227} of \num{7289} Implicit Columns flags (\qty{71.7}{\percent}). The categories cover \qty{94.1}{\percent} of class flags; \num{429} flags (\qty{5.9}{\percent}) are other or uncategorised.

Calls matching \texttt{TABLE.fields()} and \texttt{DSL.asterisk()} account for \qty{9.1}{\percent} and \qty{3.8}{\percent} of class flags. Generated DAO methods and plain SQL wildcard fragments each account for about \qty{2}{\percent}.

\begin{longtable}[hp]{|p{8cm}|S[table-format=4.0, table-column-width=3cm]|S[table-format=2.1, table-column-width=3cm]|}
	\caption{Distribution of Implicit Columns detector flags across post hoc manifestation categories.}
	\label{tab:implicitColumnsCauses}\\ \hline
	{\thead[l]{Manifestation category}} & {\thead{Detector-flag\\count}} & {\thead{Share of\\class flags}} \\
	\hline
	\endfirsthead
	\multicolumn{3}{l}
	{\tablename\ \thetable\ -- \textit{Continues...}} \\
	\hline
	{\thead[l]{Manifestation category}} & {\thead{Detector-flag\\count}} & {\thead{Share of\\class flags}} \\
	\hline
	\endhead
	\hline \multicolumn{3}{l}{\textit{Continues...}} \\
	\endfoot
	\hline
	\endlastfoot
    \texttt{DSL/DSLContext.selectFrom(TABLE)} & 3417 & 46.9 \\ \hline
    \texttt{DSL/DSLContext.select().from(TABLE)} & 1810 & 24.8 \\ \hline
    \texttt{TABLE.fields()} & 665 & 9.1 \\ \hline
    \texttt{DSL.asterisk()} & 280 & 3.8 \\ \hline
    Generated DAO methods (combined) & 163 & 2.2 \\ \hline
    \texttt{DSL/DSLContext.fetch(TABLE, ...)} & 151 & 2.1 \\ \hline
    Plain SQL wildcards & 137 & 1.9 \\ \hline
    \texttt{DSL/DSLContext.fetchOne(TABLE, ...)} & 130 & 1.8 \\ \hline
    \texttt{DSL/DSLContext.select(TABLE).from(TABLE)} & 107 & 1.5 \\ \hline
    Other/uncategorised & 429 & 5.9 \\
\end{longtable}

\subsection{Poor Man's Search Engine}

As shown in Table~\ref{tab:poorMansSearchEngineCauses}, \texttt{Field.like} accounts for \num{273} of \num{583} Poor Man's Search Engine flags (\qty{46.8}{\percent}). The four largest categories contain \num{517} flags (\qty{88.7}{\percent}). The categories cover \qty{93.7}{\percent} of class flags; \num{37} flags (\qty{6.3}{\percent}) are other or uncategorised.

The categories matching \texttt{Field.containsIgnoreCase}, \texttt{Field.likeIgnoreCase}, and \texttt{Field.contains} account for \qty{16.5}{\percent}, \qty{14.9}{\percent}, and \qty{10.5}{\percent} of class flags. Plain SQL \texttt{LIKE} fragments account for \qty{2.9}{\percent}, while \texttt{Field.likeRegex} accounts for \qty{2.1}{\percent}.

\begin{longtable}[hp]{|p{8cm}|S[table-format=3.0, table-column-width=3cm]|S[table-format=2.1, table-column-width=3cm]|}
	\caption{Distribution of Poor Man's Search Engine detector flags across post hoc manifestation categories.}
	\label{tab:poorMansSearchEngineCauses}\\ \hline
	{\thead[l]{Manifestation category}} & {\thead{Detector-flag\\count}} & {\thead{Share of\\class flags}} \\
	\hline
	\endfirsthead
	\multicolumn{3}{l}
	{\tablename\ \thetable\ -- \textit{Continues...}} \\
	\hline
	{\thead[l]{Manifestation category}} & {\thead{Detector-flag\\count}} & {\thead{Share of\\class flags}} \\
	\hline
	\endhead
	\hline \multicolumn{3}{l}{\textit{Continues...}} \\
	\endfoot
	\hline
	\endlastfoot
    \texttt{Field.like} & 273 & 46.8 \\ \hline
    \texttt{Field.containsIgnoreCase} & 96 & 16.5 \\ \hline
    \texttt{Field.likeIgnoreCase} & 87 & 14.9 \\ \hline
    \texttt{Field.contains} & 61 & 10.5 \\ \hline
    Plain SQL \texttt{LIKE} & 17 & 2.9 \\ \hline
    \texttt{Field.likeRegex} & 12 & 2.1 \\ \hline
    Other/uncategorised & 37 & 6.3 \\
\end{longtable}
